\section*{Примеры}
\addcontentsline{toc}{section}{Примеры}

\subsection*{Пример 1. Смена базиса билинейной формы}
\textbf{Идея алгоритма.} Даны матрица $B$ в старом базисе и векторы нового базиса, выраженные через старый. Координаты новых векторов записываются в столбцы матрицы перехода $C$. Новая матрица вычисляется строго по формуле тензорного закона преобразования: $\hat{B} = C^T B C$.

\subsection*{Пример 2. Выделение симметрической части квадратичной формы}
\textbf{Идея алгоритма.} Если по условию матрица квадратичной формы $B$ не является симметричной, перед любыми манипуляциями (поиск собственных значений, критерий Сильвестра и т.д.) необходимо перейти к её симметрической матрице по формуле $A = \frac{1}{2}(B + B^T)$. Только $A$ корректно отражает свойства формы.

\subsection*{Пример 3. Ортогональная диагонализация квадратичной формы}
\textbf{Идея алгоритма.} Составляется симметрическая матрица формы. Находится её характеристический многочлен и собственные значения. Диагональный вид формы будет состоять из найденных собственных значений. Для нахождения линейного преобразования ищутся собственные векторы, ортогонализуются (если кратные собственные значения) и нормируются. Из них составляется ортогональная матрица перехода $C$.

\subsection*{Пример 4. Критерий Сильвестра с параметром}
\textbf{Идея алгоритма.} Выписывается симметрическая матрица $A(\alpha)$, зависящая от параметра. Вычисляются последовательно миноры $\Delta_1, \Delta_2, \Delta_3$. Для положительной определенности решается система неравенств $\{\Delta_i > 0\}$. Для отрицательной — система $\{\Delta_1 < 0, \Delta_2 > 0, \dots\}$.
